\chapter{Hives}

\section*{Definition and Overview}
Hives, also known as urticaria, is a skin reaction in which itchy, raised, red or skin-coloured welts (wheals) appear on the skin. They occur when immune cells in the skin release histamine, causing small blood vessels to leak fluid into the skin. Hives can be triggered by foods, medicines, infections, insect stings, physical factors such as heat or pressure, and sometimes by no identifiable cause. Individual welts usually fade within 24 hours, though new ones can appear for days or weeks. Most cases of hives are short-lived and harmless, but hives with breathing difficulty or collapse indicate a severe allergic reaction (anaphylaxis) and are an emergency.

\section*{Epidemiology}
Hives are extremely common: around 15--25\% of people experience them at some point in life. Acute hives, lasting less than six weeks, are most common in children and young adults and are often linked to infections. Chronic urticaria, lasting more than six weeks, affects around 1--2\% of the population and is more common in women and in people aged 20--40. Anaphylaxis is much rarer, and food and insect triggers are the leading causes.

\section*{Aetiology and Pathophysiology}
\begin{itemize}
    \item \textbf{Foods:} peanuts, tree nuts, shellfish, eggs, and milk are common triggers
    \item \textbf{Medicines:} antibiotics, aspirin, and other drugs
    \item \textbf{Infections:} viral infections are the most common cause in children
    \item \textbf{Physical factors:} heat, cold, pressure, sunlight, exercise, and sweating cause physical urticaria
    \item \textbf{Insect stings and bites:} can cause local or widespread hives
    \item \textbf{Chronic urticaria:} often no trigger is found; an autoimmune mechanism is suspected
    \item \textbf{Pathophysiology:} mast cells release histamine, which dilates small vessels and leaks fluid, raising welts; in anaphylaxis, the reaction is systemic and life-threatening
\end{itemize}

\section*{Clinical Features}
\begin{itemize}
    \item Raised, itchy welts ranging from a few millimetres to large patches
    \item Welts may be red, pale in the centre, or skin-coloured
    \item They change shape, spread, and move around the body
    \item Individual welts usually resolve within 24 hours
    \item Angioedema: deeper swelling of the lips, eyelids, hands, feet, or genitals
    \item Symptoms can recur for weeks (acute) or months (chronic)
    \item Severe reactions add difficulty breathing, throat tightness, dizziness, and collapse
\end{itemize}

\section*{Differential Diagnosis}
\begin{itemize}
    \item Insect bites, which persist as individual lesions longer than hives
    \item Eczema and contact dermatitis, which are scaly or weeping
    \item Drug eruptions and other rashes
    \item Urticarial vasculitis, in which welts last more than 24 hours and may leave marks
    \item Mastocytosis, a rare condition of excess mast cells
    \item The transient, migrating nature of wheals is characteristic of hives
\end{itemize}

\section*{When to Seek Medical Care}
See a doctor for hives that persist beyond a few days, are severe, or have no obvious cause, and for recurrent episodes to identify triggers. Any episode with breathing difficulty or collapse is an emergency.

\textbf{Call 000 (or local emergency services) for:}
\begin{itemize}
    \item Difficulty breathing, wheezing, or throat tightness with hives
    \item Swelling of the tongue, lips, or face with breathing problems
    \item Dizziness, fainting, or collapse (anaphylaxis)
    \item Hives after a known severe allergen with any systemic symptoms
    \item Vomiting, abdominal pain, or diarrhoea with hives
    \item Use an adrenaline autoinjector immediately if available
\end{itemize}

\section*{Diagnosis and Investigations}
\begin{itemize}
    \item \textbf{History:} timing, triggers, foods, medicines, infections, and past reactions
    \item \textbf{Examination:} confirms the typical wheals and any angioedema
    \item \textbf{Trigger diary:} helps identify patterns in chronic urticaria
    \item \textbf{Allergy testing:} skin or blood tests for suspected specific allergens
    \item \textbf{Physical challenge testing:} for physical urticaria
    \item \textbf{Blood tests:} for chronic urticaria, including autoimmune markers and thyroid tests
    \item \textbf{Assessment of chronic cases:} referral to an allergy specialist where appropriate
\end{itemize}

\section*{Management}
\begin{itemize}
    \item \textbf{Antihistamines:} non-drowsy antihistamines are the first-line treatment for symptom relief
    \item \textbf{Avoid triggers:} identified foods, medicines, and physical triggers
    \item \textbf{Cool compresses and loose clothing:} relieve itching and discomfort
    \item \textbf{Chronic urticaria:} higher-dose antihistamines, and additional medicines such as montelukast or omalizumab for resistant cases
    \item \textbf{Anaphylaxis management:} adrenaline autoinjector for people with severe allergies, plus an action plan
    \item \textbf{Treat underlying infection:} when relevant
    \item \textbf{Follow-up:} chronic urticaria is managed until it settles, often within months to a few years
\end{itemize}

\section*{Complications}
\begin{itemize}
    \item Anaphylaxis in severe reactions, which is life-threatening
    \item Angioedema of the airway, causing breathing obstruction
    \item Chronic urticaria causing persistent itch, sleep disturbance, and reduced quality of life
    \item Scarring and skin changes from severe scratching
    \item Anxiety and avoidance of normal activities
    \item Interference with work, school, and sport
\end{itemize}

\section*{Prognosis and Outcomes}
Most acute hives settle within days, and antihistamines control symptoms effectively. Chronic urticaria resolves in the majority of people within 1--5 years, though it can recur. For people with severe allergies, adrenaline autoinjectors and avoidance plans prevent life-threatening outcomes. Overall, hives are a common, treatable condition, and severe complications are preventable with prompt recognition and treatment.

\section*{Prevention and Self-Care}
\begin{itemize}
    \item Identify and avoid known triggers
    \item Keep an antihistamine available for episodes
    \item Use cool compresses and loose, comfortable clothing
    \item Avoid hot showers, tight clothing, and scratching
    \item Keep a symptom and food diary for chronic hives
    \item People with severe allergies carry adrenaline and an action plan
    \item Seek urgent care for any breathing difficulty or collapse
    \item See a doctor for persistent or unexplained hives
\end{itemize}

\section*{Sources and Further Reading}
The following reputable sources provide further information for healthcare students and patients seeking reliable, current advice about hives.

\begin{itemize}
    \item Healthdirect. \textit{Hives (urticaria): causes and treatment.} https://www.healthdirect.gov.au/hives
    \item Allergy and Anaphylaxis Australia. \textit{Urticaria and allergic reactions.} https://allergyfacts.org.au/
    \item ASCIA. \textit{Chronic urticaria information.} https://www.allergy.org.au/
    \item NHS. \textit{Hives: overview and treatment.} https://www.nhs.uk/conditions/hives/
    \item Mayo Clinic. \textit{Hives and angioedema: symptoms and causes.} https://www.mayoclinic.org/
    \item MSD Manual. \textit{Urticaria.} https://www.msdmanuals.com/
\end{itemize}
